\section{Topic allocation CRDT}\label{sec:session_crdt}

The session layer maintains the topic-to-subject-ID mapping through a conflict-free
replicated data type (CRDT) shared among all nodes by periodic gossip.
Under the assumptions listed below, every node's local copy of the mapping converges
to the same collision-free, non-divergent configuration within a bounded time after
the last state-changing event.

This section is normative for the arbitration semantics of the mapping.
The primary source of truth is the TLA$^+$ model in \texttt{model/tlaplus/Core.tla}
of the Cy reference implementation; the operators and theorems described here are
a faithful transcription of that model, and the C reference implementation has been
shown to match it under the stated assumptions.

\subsection{Topic record}

Each topic the node is aware of is represented by a triple
\[
    t = (\mathit{hash},\ \mathit{evictions},\ \mathit{age})
\]
where:

\begin{description}
    \item[\textit{hash}] is the 64-bit topic hash defined in
    section~\ref{sec:session_rapidhash}.
    \item[\textit{evictions}] is a 32-bit unsigned integer incremented each time the
    topic loses arbitration on its current candidate subject-ID. It serves as a
    Lamport clock. The reserved range at the top of the counter encodes pinned
    topics as described in section~\ref{sec:session_subject_id_mapping}.
    \item[\textit{age}] is a local real or logical clock counting roughly seconds
    since the topic was first seen on the network. It is not transmitted directly
    on the wire; gossip instead carries its base-2 logarithm (see below).
\end{description}

The subject-ID of a topic is the pure function
$\mathit{subject\_id}(\mathit{hash}, \mathit{evictions})$ defined in
section~\ref{sec:session_subject_id_mapping}.

\subsection{Log-age}

CRDT operators use the quantity
\[
    \mathit{lage}(t) = \lfloor \log_2(\mathit{age\_seconds}(t)) \rfloor,
    \qquad \mathit{lage}(t) = -1 \text{ when } \mathit{age\_seconds}(t) = 0
\]
in place of the raw age.
Taking the floor of the logarithm slows the arbitration-relevant clock as topics
grow older, so that the CRDT is stable on long timescales even though the underlying
age counter never stops advancing. The encoded range is $[-1,\ 35]$, where
$\mathit{lage} = 35$ corresponds to approximately $2^{35}$ seconds, or over a
thousand years.

\begin{remark}
    The naive approach of using the raw age in arbitration would make any two
    quiescent nodes disagree on priorities every time a clock edge crosses an
    integer second boundary, which breaks the CRDT's eventual consistency.
    Reducing the comparison quantity to $\lfloor \log_2(\mathit{age}) \rfloor$
    makes the set of points where priority can change sparse on quiescent states,
    satisfying the stability assumption A9 below.
\end{remark}

Gossip messages carry a topic's $\mathit{lage}$ value rather than its age.
On reception, a node merges the incoming $\mathit{lage}$ into its local age by
adjusting the local age field so that
\[
    \mathit{local\_age} \leftarrow \max\left(\mathit{local\_age},\ 2^{\mathit{remote\_lage}}\right).
\]
This is a monotonic, idempotent join: repeated application is a no-op, and the
result is independent of the order in which gossip messages are processed.

\subsection{Arbitration}

Two topics are in \emph{collision} when they have distinct hashes but map to the
same subject-ID. Two records with the same hash but different eviction counters
are in \emph{divergence}; these are different opinions about the same logical topic
that must be reconciled.

\begin{description}
    \item[Collision arbitration] Older $\mathit{lage}$ wins. On a tie, the record
    with the numerically lower hash wins.
    \item[Divergence arbitration] Older $\mathit{lage}$ wins. On a tie, the record
    with the higher eviction count wins, because a higher counter means the record
    has observed more prior conflicts and is therefore more up to date.
\end{description}

These two predicates are defined as follows, verbatim from
\texttt{model/tlaplus/Core.tla}. \texttt{Log2Floor(x)} is
$\lfloor \log_2 x \rfloor$ with the special case $\mathrm{Log2Floor}(0) = -1$;
\texttt{\#} is the TLA\textsuperscript{+} inequality operator, $\mathtt{=}$ is
equality, $\mathtt{/\char`\\}$ is logical AND, and $\mathtt{\char`\\/}$ is
logical OR.

\begin{minted}{text}
LeftWinsCollision(a, b) == IF   Log2Floor(a.age) # Log2Floor(b.age)
                           THEN Log2Floor(a.age) > Log2Floor(b.age)
                           ELSE a.hash < b.hash

LeftWinsDivergence(a, b) == \/  Log2Floor(a.age) > Log2Floor(b.age)
                            \/ (Log2Floor(a.age) = Log2Floor(b.age) /\ a.evictions > b.evictions)
\end{minted}

\subsection{AllocateTopic operator}

\texttt{AllocateTopic} inserts a topic into the node's local topic set, resolving
any resulting collision by cascading evictions through the hash space.
The operator is recursive; termination follows from the bounded-escape property in
assumption A3 below. The TLA\textsuperscript{+} definition below is verbatim
from \texttt{model/tlaplus/Core.tla} and is reproduced in full in
appendix~\ref{sec:appendix_core_tla}.

\begin{minted}{text}
RECURSIVE AllocateTopic(_, _)
AllocateTopic(t, topics) ==
    LET ts == RemoveTopic(t.hash, topics)
        x == GetBySubjectID(SubjectID(t.hash, t.evictions), ts)
        Evicted(z) == [hash |-> z.hash, evictions |-> 1 + z.evictions, age |-> z.age]
    IN   IF x = Nothing             THEN {t} \cup ts
    ELSE IF LeftWinsCollision(t, x) THEN AllocateTopic(Evicted(x), {t} \cup ts)
    ELSE                                 AllocateTopic(Evicted(t), ts)
\end{minted}

The operator is invoked whenever a new local topic is created, whenever a gossip
forces the local record for a topic to adopt new eviction values, and whenever
a cascade from another topic displaces the current record.

\subsection{AcceptGossip operator}

\texttt{AcceptGossip} applies a received gossip message to the node's local topic
set. The operator distinguishes two cases: the gossiped topic is known locally
(its hash matches a local record) or it is not. The TLA\textsuperscript{+}
definitions below are verbatim from \texttt{model/tlaplus/Core.tla}.
\texttt{Nothing} denotes the absence of a matching record;
\texttt{FloorToPow2(x)} returns $2^{\lfloor \log_2 x \rfloor}$ with the
special case $\mathrm{FloorToPow2}(0) = 0$.

\begin{minted}{text}
AcceptGossip(remote, topics) ==
    IF GetByHash(remote.hash, topics) # Nothing     \* If the topic exists locally
    THEN AcceptGossip_Divergence(remote, topics)    \* Divergence resolver ensures collision-freedom
    ELSE AcceptGossip_Collision(remote, topics)     \* We don't know this topic, just ensure collision-freedom

AcceptGossip_Divergence(remote, topics) ==
    LET hash == remote.hash
        local == GetByHash(hash, topics)
        \* The serialization protocol floors the remote age to pow2 before transmission.
        new_age == Max({IF local # Nothing THEN local.age ELSE 0, FloorToPow2(remote.age)})
    IN
    IF local # Nothing THEN
        IF local.evictions # remote.evictions /\ LeftWinsDivergence(remote, local)
        THEN AllocateTopic([hash|->hash, evictions|->remote.evictions, age|->new_age], topics)
        ELSE              {[hash|->hash, evictions|-> local.evictions, age|->new_age]} \cup RemoveTopic(hash, topics)
    ELSE topics

AcceptGossip_Collision(remote, topics) ==
    LET local == GetBySubjectID(SubjectID(remote.hash, remote.evictions), topics)
    IN IF local # Nothing /\ LeftWinsCollision(remote, local)
    THEN AllocateTopic([hash |-> local.hash, evictions |-> local.evictions + 1, age |-> local.age], topics)
    ELSE topics
\end{minted}

\begin{remark}[breakable]
    A naive implementation could always run the collision resolver after the
    divergence resolver, i.e.\ compute
    \texttt{AcceptGossip\_Collision(remote, AcceptGossip\_Divergence(remote, topics))}.
    This is correct but suboptimal: the local record may already win the
    divergence arbitration against the remote, in which case the remote record
    should be ignored entirely; running the collision resolver afterwards can make
    a local record lose a collision to a stale remote record and trigger
    unnecessary eviction-counter increments. The divergence-first rule above
    avoids that outcome -- when the local record wins the divergence arbitration,
    the collision pass is skipped because divergence resolution implicitly keeps
    the local mapping collision-free.
\end{remark}

\subsection{Pinned topics}

Pinned topics are ordinary CRDT records with their eviction counter pre-set to
$\text{UINT32\_MAX} - \text{pinned\_subject\_id}$, as defined in
section~\ref{sec:session_subject_id_mapping}.
They participate in arbitration on the same terms as non-pinned topics: an older
non-pinned topic can displace a younger pinned newcomer, and an older pinned topic
can displace a younger conflicting pinned topic.
Pinning is therefore a \emph{request}, not a guarantee.

\subsection{Assumptions}

The convergence guarantee of the allocation protocol rests on the following
assumptions. Some are structural properties of the mapping function and are
satisfied by construction; others are deployment constraints that compliant
implementations and system integrators are required to uphold. The numbering is
consistent with the accompanying proof in \texttt{model/proof.md}.

\begin{description}
    \item[A1. Common deterministic mapping] All nodes use the same subject-ID
    mapping function and the same parameters -- in particular, the same modulus.
    \item[A2. Hash identity] Active topic names are hash-unique. The probability of
    a 64-bit hash collision is negligible and is ignored by the protocol.
    \item[A3. Bounded escape] For any hash, any eviction value, and any blocked
    subject-ID set smaller than the capacity bound $K$, there exists a small
    increment within $K$ steps that places the topic at a subject-ID outside the
    blocked set. For quadratic probing on a prime modulus $M$ with $M \bmod 4 = 3$
    this holds with $K \approx \lfloor M / 2 \rfloor + 1$.
    \item[A4. Capacity bound] The number of active non-pinned topics $N$ satisfies
    $N \le K$.
    \item[A5. Gossip fairness] Every live topic record is eventually received and
    processed by every relevant node infinitely often.
    \item[A6. Progress fairness] Each node keeps executing the protocol; resource
    exhaustion or scheduling starvation does not halt a live node.
    \item[A7. Resource sufficiency] Each node has enough memory and CPU to process
    and complete the local cascade triggered by any single allocation step.
    \item[A8. Eviction counter safety] Eviction counters do not overflow in
    reachable executions.
    \item[A9. Stability window] Stable windows longer than the worst-case repair
    time $T_{\mathrm{repair}}$ exist, during which the active topic set and
    pairwise arbitration priorities do not flip.
\end{description}

\subsection{Convergence}

Under assumptions A1--A9, any execution restricted to a single stability window
converges in finite time to a state that is free of collisions and divergences.
The argument is given in full in \texttt{model/proof.md}; the key idea is that
higher-priority topics are never displaced by lower-priority topics, so the
arbitration induces a fixed total order during the window and each topic
stabilizes by induction on its priority rank.

When the stability window ends -- for example, because a long-lived topic's
$\mathit{lage}$ steps up and flips a tie -- some transient churn may briefly
reappear. Assumption A9 guarantees that the next stable window is long enough to
repair the churn before the next flip, and the overall execution remains eventually
consistent.

\subsection{Capacity}

For a quadratic probing sequence on a prime modulus $M$ with $M \bmod 4 = 3$ the
collision-free probing capacity is approximately $\lfloor M / 2 \rfloor$
(assumption A3). The number of active non-pinned topics should not exceed this
bound; above it, the arbitration behaviour remains well-defined but convergence
within a single stability window is no longer guaranteed.
